\documentclass[10pt]{article}

\usepackage[margin=1in]{geometry}
\usepackage{amsmath}
\usepackage{amssymb}
\usepackage{amsthm}
\usepackage{mathtools}
\usepackage[shortlabels]{enumitem}
\usepackage[normalem]{ulem}
\usepackage{courier}

\usepackage{hyperref}
\hypersetup{
  colorlinks   = true, %Colours links instead of ugly boxes
  urlcolor     = black, %Colour for external hyperlinks
  linkcolor    = blue, %Colour of internal links
  citecolor    = blue  %Colour of citations
}

\usepackage[T1]{fontenc}
\usepackage{listings}
\lstset{
    language=HTML
    ,basicstyle=\linespread{1}\ttfamily
    ,keywordstyle=
    %,numbers=left
    ,breaklines=true
    }

%%%%%%%%%%%%%%%%%%%%%%%%%%%%%%%%%%%%%%%%%%%%%%%%%%%%%%%%%%%%%%%%%%%%%%%%%%%%%%%%

\theoremstyle{definition}
\newtheorem{problem}{Problem}
\newtheorem{note}{Note}
\newcommand{\E}{\mathbb E}
\newcommand{\R}{\mathbb R}
\DeclareMathOperator{\Var}{Var}
\DeclareMathOperator*{\argmin}{arg\,min}
\DeclareMathOperator*{\argmax}{arg\,max}

\newcommand{\trans}[1]{{#1}^{T}}
\newcommand{\loss}{\ell}
\newcommand{\w}{\mathbf w}
\newcommand{\mle}[1]{\hat{#1}_{\textit{mle}}}
\newcommand{\map}[1]{\hat{#1}_{\textit{map}}}
\newcommand{\normal}{\mathcal{N}}
\newcommand{\x}{\mathbf x}
\newcommand{\y}{\mathbf y}
\newcommand{\ltwo}[1]{\lVert {#1} \rVert}

%%%%%%%%%%%%%%%%%%%%%%%%%%%%%%%%%%%%%%%%%%%%%%%%%%%%%%%%%%%%%%%%%%%%%%%%%%%%%%%%

\begin{document}
\begin{center}
    {
\Large
    Practice Problems: Counting
}

    \vspace{0.1in}
\end{center}

\section{Review Problems}

\begin{note}
    Your quiz will have python review problems, but they will be in a different format.
    Instead of being presented as "raw" python,
    they will be embedded with the shell.
\end{note}

\filbreak
\begin{problem}
    Write the output of the final command in the following terminal session.
    If the command has no output, then leave the problem blank.
\end{problem}
\begin{lstlisting}
$ cd; rm -rf quiz; mkdir quiz; cd quiz
$ cat > foo.py <<EOF
xs = [1, 2, 3]
def foo(y):
    xs.append(y)
    return y
foo(5)
print('sum(xs)=', sum(xs))
EOF
$ python3 foo.py
\end{lstlisting}

\section{Nested for loops}

\begin{note}
    For the rest of these problems, the goal is to learn shortcuts to counting the number of times certain lines will get run.
    For nested for loops, we multiply the number of iterations of the loop;
    for sequential loops, we add.
\end{note}

\filbreak
\begin{problem}
    Write the output of the final command in the following terminal session.
    If the command has no output, then leave the problem blank.
\end{problem}
\begin{lstlisting}
$ cd; rm -rf quiz; mkdir quiz; cd quiz
$ N=8
$ cat > foo.py <<EOF
count = 0
for i in range($N):
    count += 1
print('count=', count)
EOF
$ python3 foo.py
\end{lstlisting}
\vspace{0.4in}

\filbreak
\begin{problem}
    Write the output of the final command in the following terminal session.
    If the command has no output, then leave the problem blank.
\end{problem}
\begin{lstlisting}
$ cd; rm -rf quiz; mkdir quiz; cd quiz
$ N=8
$ cat > foo.py <<EOF
count = 0
for i in range($N):
    for j in range($N):
        count += 1
print('count=', count)
EOF
$ python3 foo.py
\end{lstlisting}
\vspace{0.4in}

\filbreak
\begin{problem}
    Write the output of the final command in the following terminal session.
    If the command has no output, then leave the problem blank.
\end{problem}
\begin{lstlisting}
$ cd; rm -rf quiz; mkdir quiz; cd quiz
$ N=8
$ cat > foo.py <<EOF
count = 0
for i in range($N):
    for j in range($N):
        for k in range($N):
            count += 1
print('count=', count)
EOF
$ python3 foo.py
\end{lstlisting}
\vspace{0.1in}

\filbreak
\section{Sequential for loops}

\begin{problem}
    Write the output of the final command in the following terminal session.
    If the command has no output, then leave the problem blank.
\end{problem}
\begin{lstlisting}
$ cd; rm -rf quiz; mkdir quiz; cd quiz
$ N=8
$ cat > foo.py <<EOF
count = 0
for i in range($N):
    count += 1
for i in range($N):
    count += 1
print('count=', count)
EOF
$ python3 foo.py
\end{lstlisting}
\vspace{0.1in}

\filbreak
\begin{problem}
    Write the output of the final command in the following terminal session.
    If the command has no output, then leave the problem blank.
\end{problem}
\begin{lstlisting}
$ cd; rm -rf quiz; mkdir quiz; cd quiz
$ N=8
$ cat > foo.py <<EOF
count = 0
for i in range($N):
    count += 1
for i in range($N):
    count += 1
for i in range($N):
    count += 1
print('count=', count)
EOF
$ python3 foo.py
\end{lstlisting}
\vspace{0.4in}

\filbreak
\section{Combinations of sequential and nested loops}

\begin{problem}
    Write the output of the final command in the following terminal session.
    If the command has no output, then leave the problem blank.
\end{problem}
\begin{lstlisting}
$ cd; rm -rf quiz; mkdir quiz; cd quiz
$ N=8
$ cat > foo.py <<EOF
count = 0
for i in range($N):
    for j in range($N):
        count += 1
for i in range($N):
    count += 1
    count += 1
count += 1
print('count=', count)
EOF
$ python3 foo.py
\end{lstlisting}
\vspace{0.4in}

\filbreak
\begin{problem}
    Write the output of the final command in the following terminal session.
    If the command has no output, then leave the problem blank.
\end{problem}
\begin{lstlisting}
$ cd; rm -rf quiz; mkdir quiz; cd quiz
$ N=8
$ cat > foo.py <<EOF
count = 0
for i in range($N):
    count += 1
    for j in range($N):
        count += 1
for i in range($N):
    count += 1
print('count=', count)
EOF
$ python3 foo.py
\end{lstlisting}
\vspace{0.4in}

\filbreak
\begin{problem}
    Write the output of the final command in the following terminal session.
    If the command has no output, then leave the problem blank.
\end{problem}
\begin{lstlisting}
$ cd; rm -rf quiz; mkdir quiz; cd quiz
$ N=8
$ cat > foo.py <<EOF
count = 0
for i in range($N):
    count += 1
for i in range($N):
    count += 1
    for j in range($N):
        count += 1
    count += 1
    for j in range($N):
        count += 1
        for k in range($N):
            count += 1
for i in range($N):
    count += 1
print('count=', count)
EOF
$ python3 foo.py
\end{lstlisting}
\vspace{0.4in}

\filbreak
\section{``Simple'' Expressions in Multiple Variables}

\begin{problem}
    Write the output of the final command in the following terminal session.
    If the command has no output, then leave the problem blank.
\end{problem}
\begin{lstlisting}
$ cd; rm -rf quiz; mkdir quiz; cd quiz
$ N=8
$ M=4
$ cat > foo.py <<EOF
count = 0
for i in range($N):
    for j in range($M):
        count += 1
print('count=', count)
EOF
$ python3 foo.py
\end{lstlisting}
\vspace{0.4in}

\filbreak
\begin{problem}
    Write the output of the final command in the following terminal session.
    If the command has no output, then leave the problem blank.
\end{problem}
\begin{lstlisting}
$ cd; rm -rf quiz; mkdir quiz; cd quiz
$ N=8
$ M=4
$ cat > foo.py <<EOF
count = 0
for i in range($N):
    count += 1
for j in range($M):
    count += 1
print('count=', count)
EOF
$ python3 foo.py
\end{lstlisting}
\vspace{0.4in}

\filbreak
\begin{problem}
    Write the output of the final command in the following terminal session.
    If the command has no output, then leave the problem blank.
\end{problem}
\begin{lstlisting}
$ cd; rm -rf quiz; mkdir quiz; cd quiz
$ N=8
$ M=4
$ cat > foo.py <<EOF
count = 0
for i in range($N):
    count += 1
    for j in range($M):
        count += 1
    for j in range($N):
        count += 1
for i in range($M):
    count += 1
    for j in range($M):
        count += 1
        count += 1
        count += 1
print('count=', count)
EOF
$ python3 foo.py
\end{lstlisting}
\vspace{0.4in}

\filbreak
\begin{problem}
    Write the output of the final command in the following terminal session.
    If the command has no output, then leave the problem blank.
\end{problem}
\begin{lstlisting}
$ cd; rm -rf quiz; mkdir quiz; cd quiz
$ N=8
$ M=4
$ O=2
$ cat > foo.py <<EOF
count = 0
for i in range($N):
    count += 1
    for j in range($M):
        count += 1
    for j in range($O):
        count += 1
for i in range($M):
    count += 1
    for j in range($O):
        count += 1
        count += 1
        count += 1
print('count=', count)
EOF
$ python3 foo.py
\end{lstlisting}
\vspace{0.4in}

\filbreak
\section{``Complex'' Expressions in Multiple Variables}

\begin{problem}
    Write the output of the final command in the following terminal session.
    If the command has no output, then leave the problem blank.
\end{problem}
\begin{lstlisting}
$ cd; rm -rf quiz; mkdir quiz; cd quiz
$ N=8
$ M=4
$ O=2
$ cat > foo.py <<EOF
count = 0
for i in range($N * $M):
    for j in range($M * $O):
        count += 1
print('count=', count)
EOF
$ python3 foo.py
\end{lstlisting}
\vspace{0.4in}

\filbreak
\begin{problem}
    Write the output of the final command in the following terminal session.
    If the command has no output, then leave the problem blank.
\end{problem}
\begin{lstlisting}
$ cd; rm -rf quiz; mkdir quiz; cd quiz
$ N=8
$ M=4
$ O=2
$ cat > foo.py <<EOF
count = 0
for i in range($N + $M):
    for j in range($M * $O):
        count += 1
print('count=', count)
EOF
$ python3 foo.py
\end{lstlisting}
\vspace{0.4in}

\filbreak
\begin{problem}
    Write the output of the final command in the following terminal session.
    If the command has no output, then leave the problem blank.
\end{problem}
\begin{lstlisting}
$ cd; rm -rf quiz; mkdir quiz; cd quiz
$ N=8
$ M=4
$ O=2
$ cat > foo.py <<EOF
count = 0
for i in range($N + $M):
    count += 1
for i in range($M * $O):
    count += 1
print('count=', count)
EOF
$ python3 foo.py
\end{lstlisting}
\vspace{0.4in}

\filbreak
\begin{problem}
    Write the output of the final command in the following terminal session.
    If the command has no output, then leave the problem blank.
\end{problem}
\begin{lstlisting}
$ cd; rm -rf quiz; mkdir quiz; cd quiz
$ N=8
$ M=4
$ O=2
$ cat > foo.py <<EOF
count = 0
for i in range($N + $M):
    count += 1
    for j in range($N ** $O):
        count += 1
for i in range($M * $O):
    count += 1
print('count=', count)
EOF
$ python3 foo.py
\end{lstlisting}
\vspace{0.4in}

\filbreak
\begin{problem}
    Write the output of the final command in the following terminal session.
    If the command has no output, then leave the problem blank.
\end{problem}
\begin{lstlisting}
$ cd; rm -rf quiz; mkdir quiz; cd quiz
$ N=8
$ M=4
$ O=2
$ cat > foo.py <<EOF
count = 0
for i in range($N + $M):
    count += 1
for i in range($N ** $O):
    count += 1
    for j in range($M + $O):
        count += 1
print('count=', count)
EOF
$ python3 foo.py
\end{lstlisting}
\vspace{0.4in}

\filbreak

\begin{note}
    The only math operations I will use in the real quiz are addition, multiplication, and exponentiation.
\end{note}

\end{document}
